\par
(2.15) 上述讨论的一部分可推广到高维。设 $S$ 是维数为 $n$ 的正则局部环的谱，闭点为 $s$，且 $k(s)$ 是 $k$ 的有限扩张。仍可考虑从 $S$ 逐次爆破闭点所得的概形 $f:S'\longrightarrow S$ 的集合 $I$，并如 2.1 一样，置

\begin{equation*}
S^*=\varinjlim_{S'\in I}\operatorname{Irr}\bigl(f^{-1}(s)\bigr).
\end{equation*}
集合 $S^*$ 仍称为\underline{$S$ 中与 $s$ 无限邻近的点}的集合。

若 $X$ 是 $S$ 的闭子概形，且 $t\in S^*$，则如 2.8 一样定义 $X$ 在 $t$ 处的重数。

设 $f:S'\longrightarrow S$ 是若干闭点处爆破态射的复合，$D$ 是 $S$ 上的除子，$D^*=f^*D$ 是其在 $S'$ 上的全变换。

%%%%%%%%%% source scan idx 34 | folio 27 %%%%%%%%%%

设 $E$ 是 $S$ 上的曲线，$E'$ 是其在 $S'$ 上的严格变换，假设 $D\cap E$ 的支集包含于 $s$。投影公式仍给出

\begin{equation}\tag{2.16}
(D,E)=(D^*,E').
\end{equation}
仿照 2.9，借助（2.16）可证明

\begin{equation}\tag{2.17}
(D,E)=\sum_{t\in S^*}\mu_t(D)\mu_t(E)[k(t):k] \quad .
\end{equation}
如 2.5 一样定义邻接点。仿照 2.11 可验证

\begin{equation}\tag{2.18}
\mu_t(E)=
\sum_{\substack{t'\ \text{邻接}\\\text{于}\ t}}
\mu_{t'}(E)[k(t'):k(t)] \quad .
\end{equation}
另一方面，公式 2.11(ii) 一般不成立（见 2.25 的例子）。

2.19. 设 $S$ 是域 $k$ 上的正则射影曲面。用 $S^*$ 表示对 $S$ 的闭点所取的不交并

\begin{equation*}
S^*=\coprod_{s\in S}\operatorname{Spec}(\underline{O}_{s,S})^*,
\end{equation*}
并称 $S^*$ 为\underline{与 $S$ 的点无限邻近的点}的集合。

设 $I$ 是从 $S$ 逐次爆破闭点所得曲面的集合。在 $S$ 的函数域的正则模型按控制关系排序的集合中，集合 $I$ 是共尾的。对 $i\in I$，用 $S_i$ 表示相应曲面，带有态射 $p_i:S_i\longrightarrow S$。用 $\operatorname{Div}^{\pm}(S_i)$ 表示 $S_i$ 上的除子群，除子有效与否均可。

\underline{定义} 2.20. (i) \underline{$S$ 上的“曲线”群是归纳极限}

\begin{equation*}
\varinjlim_{i\in I}\operatorname{Div}^{\pm}(S_i).
\end{equation*}

(ii) \underline{$S$ 上的“可逆层”范畴是归纳极限}

%%%%%%%%%% source scan idx 35 | folio 28 %%%%%%%%%%

（SGA 4 VI 6），\underline{其中取的是各 $(S_i)_{i\in I}$ 上的可逆层范畴。}

这些对象都是 $S$ 的双有理不变量。

\underline{例} 2.21. (i) $S$ 上的每条“曲线”$C$ 都定义 $S$ 上的一个“可逆层”$\underline{O}(C)$。

(ii) $S$ 上的每个除子 $C$ 都定义 $S$ 上的一条“曲线”，记为 $C^*$。

(iii) 若 $t\in S^*$，则如 2.3 一样定义曲面 $S_+(t)$ 及 $S_+(t)$ 上的除子 $t^*$；此除子所定义的“曲线”仍记为 $t^*$。

(2.22) 容易验证，一条“曲线”$C$（分别地，一个“可逆层”$\mathcal L$）可用唯一的方式写成

\begin{equation*}
C=C_0^*-\sum_{t\in S^*}i_t t^*
\qquad\text{分别地}\qquad
\mathcal L=\mathcal L_0\otimes\underline{O}
\left(-\sum_{t\in S^*}i_t t^*\right),
\end{equation*}
其中 $C_0$（分别地，$\mathcal L_0$）是 $S$ 上的除子（分别地，可逆层），且整数 $i_t$ 几乎全为零。称 $C$（分别地，$\mathcal L$）为赋予虚重数 $i_t$ 的曲线 $C_0$（分别地，可逆层 $\mathcal L_0$）。

$S$ 上的“曲线”群 $\operatorname{Div}^{\pm}(S)\times\mathbb Z^{(S^*)}$ 的完备化 $\operatorname{Div}^{\pm}(S)\times\mathbb Z^{S^*}$ 仍是 $S$ 的双有理不变量；称为 $S$ 上的“广义曲线”群。在 $S$ 上，“广义曲线”$C$ 由一个赋予了一族可能无穷多个虚重数的除子 $C_0$ 表示。

$S$ 上的“广义可逆层”以同样方式定义。

%%%%%%%%%% source scan idx 36 | folio 29 %%%%%%%%%%

\underline{例} 2.23. 假设 $k$ 完美。则对每个点 $t\in S^*$ 赋予虚重数 $-1$ 的可逆层 $\Omega^2_{S/k}$ 是 $S$ 的双有理不变量；将其记为 $\Omega^2_{S/k}{}^*$。

\underline{例} 2.24. 设 $f:S_1\longrightarrow S_2$ 是 $k$ 上正则射影曲面之间的支配态射。取 $S_2$ 的一个模型，使得 $S_2$ 上的“曲线”或“可逆层”成为通常的除子或可逆层，便可无困难地定义它的逆像。取极限（我没有核查这种取极限是否可行），也可类似地定义“广义曲线”或“广义可逆层”的逆像。假设 $k$ 完美，且 $f$ 在泛点处平展。则 $f$ 的\underline{不同除子（different）} $\mathfrak d(f)$ 是 $S_1$ 上由下式刻画的有效除子

\begin{equation*}
f^*\Omega^2_{S_2/k}\bigl(\mathfrak d(f)\bigr)
\overset{\sim}{\longrightarrow}\Omega^2_{S_1/k} \quad .
\end{equation*}
此不同除子不是 $f$ 的双有理不变量；要得到双有理不变量，必须考虑广义重合曲线

\begin{equation*}
\mathfrak d^*(f)=\mathfrak d(f)+\sum_{t\in S_1^*}t^*
-f^*\left(\sum_{t\in S_2^*}t^*\right),
\end{equation*}
或简写为

\begin{equation*}
\mathfrak d^*(f)=\Omega^2_{S_1/k}{}^*
-f^*\Omega^2_{S_2/k}{}^* \quad .
\end{equation*}
此“广义曲线”的支集包含于 $S_1$ 中 $f$ 不平展的闭子集。